One of the main implementation claims of this work is that achieving
strong onboard vision performance on a compact MCU-plus-EdgeTPU platform
depends less on raw inference speed alone and more on how efficiently
the host prepares, places, and transfers data. This is particularly
visible in the code added relative to the initial fork, where the
runtime evolves from a board example into a carefully staged real-time
system.

The core execution environment is FreeRTOS. Rather than running
perception as a single sequential loop, the implementation decomposes
the workload into cooperating tasks with explicit ownership boundaries.
The most relevant example is the continuous detection pipeline, where
the code in \texttt{detection\_task.cc} organizes the system around a
preparation task and an inference task. The preparation task acquires
the latest frame, scales and converts it through the hardware pixel
pipeline, applies quantization when needed, and writes the result into a
staging buffer. The inference task copies the prepared staging contents
into the EdgeTPU input tensor, releases the staging buffer immediately,
and then executes inference and post-processing while preparation for
the next frame can already proceed. This design is not simply an
optimization; it is the mechanism through which the implementation turns
limited memory and compute resources into sustained throughput.

An important low-level detail is that the runtime does not feed the
detector directly from the raw camera bus format. The active camera
configuration captures native 1280x720 frames into 64-byte-aligned
framebuffers with a 4-byte-per-pixel host representation. In practice,
the software treats these buffers as XRGB8888-like surfaces because this
layout is easier to move through CSI DMA, cache maintenance, and the
RT1176 pixel pipeline than a tightly packed 3-byte RGB tensor. This is
precisely where the PXP block becomes essential rather than optional: it
converts the aligned capture surface into RGB888 output at the detector
input size, eliminating expensive per-pixel software shuffles on the
Cortex-M7.

This organization also clarifies the broader software split introduced
in the architecture chapter. The lower implementation layer is
task-oriented in C++, meaning that concurrency is expressed directly
through FreeRTOS tasks, semaphores, queues, and explicit buffer handoff
between producer and consumer stages. By contrast, the upper user layer
remains script-oriented in MicroPython, where these same services are
invoked as high-level operations without exposing the user to the
internal synchronization structure. The runtime and memory subsystem is
precisely where this separation becomes most meaningful: scripts can
request model loading, camera capture, or continuous detection, but the
correctness and performance of those operations depend on the
task-oriented C++ layer maintaining deterministic ownership over memory,
timing, and accelerator access.

The use of double-buffered staging is especially important because it
enforces a clean ownership discipline between producer and consumer
stages. The staging buffer is written by the preparation side and read
by the inference side under semaphore-controlled handoff, while the TPU
tensor is written only when the staging contents are ready. This
minimizes race conditions, reduces needless synchronization complexity,
and preserves the ability to overlap work. In a larger system with
abundant memory, one could compensate for poor orchestration by brute
force. In the present implementation, orchestration is the only
realistic way to maintain high frame rate while preserving the low-power
character of the platform.

The same section of code also shows why byte alignment becomes a
first-order systems concern. The staging area in SDRAM is explicitly
aligned to cache boundaries, while the PXP conversion path performs
cache clean-and-invalidate operations before DMA writes and invalidation
after completion so that stale boundary lines do not overwrite fresh
image data. This is needed because the TFLite arena is not naturally
aligned to the same cache-line granularity as the DMA engine. Without
that discipline, a seemingly minor mismatch between 16-byte tensor
allocation and 32-byte cache-line behavior would produce intermittent
corruption exactly at the stage where the detector input is assembled.

The memory layout reflects the same philosophy. The linker script
introduces a deliberate partitioning across internal memory, OCRAM,
heap, SDRAM, and a dedicated region for camera-related traffic.
MicroPython is placed explicitly in OCRAM rather than in the most
latency-sensitive memory regions, while large buffers, tracker state,
and staging areas are pushed into SDRAM where capacity matters more than
absolute access latency. This is visible both in the custom linker
script and in the pervasive use of explicit section placement such as
\texttt{.sdram\_bss}. The implementation thus accepts that not all
memory is equal and uses the linker as part of the optimization strategy
rather than as a passive build artifact.

This memory strategy also explains a design choice that would otherwise
look surprising from a purely model-centric viewpoint: the system first
tolerates a larger, alignment-friendly framebuffer representation and
only later collapses it into the compact RGB tensor expected by the
model. In other words, the runtime pays for structure before it pays for
compactness. That tradeoff is justified because aligned DMA and
deterministic hardware conversion are cheaper on this platform than
repeatedly manipulating tightly packed pixels in software.

This issue becomes more acute once YOLO-class detectors are considered.
Larger detector inputs significantly increase the pressure on staging
buffers, copied tensors, post-processing data, and runtime metadata. In
a conventional desktop pipeline this is almost invisible. On an MCU
platform it becomes decisive. The code therefore embodies a
memory-rearrangement strategy aimed at fitting the MicroPython runtime,
the real-time tasks, the tracker structures, and the camera-to-TPU
pipeline into a workable whole while minimizing inefficient copies. This
is one of the places where the implementation differs most strongly from
a conceptual prototype: the problem is no longer simply to run a model,
but to run it without collapsing the rest of the system.

The same logic informs the build system. The custom CMake configuration
aggregates a full embedded MicroPython build, nanopb-generated protocol
code, additional audio and communication components, and the custom
runtime sources into a single executable with a purpose-built linker
script. This makes the firmware itself part of the experimental method.
The implementation is not a loose collection of demos, but a
custom-built runtime whose compilation model, memory map, and
dependencies are aligned with the targeted deployment conditions.

From a research perspective, this subsection supports a broader claim:
on compact embedded AI platforms, the practical efficiency of a
specialized accelerator depends on the host system's ability to keep the
accelerator supplied with valid input at the right pace. The
implementation therefore treats memory organization and RTOS scheduling
as intrinsic parts of accelerator utilization rather than as secondary
engineering details.
